%==============================================================================
% Capitulo 8: Discusion
%==============================================================================
\section{Discusion}

\subsection{Interpretacion de los resultados}

Los resultados obtenidos confirman la hipotesis de trabajo formulada en la Seccion~1.4. La reducción del 89~\% en el tiempo medio de registro y la exactitud global del 92~\% alcanzada por el clasificador automático superan ampliamente los umbrales planteados en los objetivos específicos del trabajo. Tres aspectos de los resultados merecen una interpretación detallada.

La reducción de tiempo se sostiene incluso al incorporar el costo de inferencia del modelo de lenguaje grande externo, dado que la primera etapa deterministica del clasificador resuelve aproximadamente el 62~\% de los casos sin invocar al modelo de Google. Este hallazgo valida la hipotesis subsidiaria sobre la eficiencia del esquema hibrido: el filtro deterministico no solo reduce el costo económico del sistema, sino que además disminuye la latencia mediana del flujo automatizado al evitar llamadas de red innecesarias al servicio externo. Desde una perspectiva operativa, esta propiedad del clasificador introduce un punto de inflexion en la relacion costo-beneficio de los sistemas basados en LLM: el costo marginal de la inferencia se aplica únicamente a los casos que realmente demandan capacidad semántica, y no al volumen completó de incidentes.

La dispersión observada en los tiempos del flujo automatizado ---con un rango de 11 a 31 segundos, frente a los 96 a 289 segundos del flujo manual--- evidencia una compresion sustancial de la variabilidad. Mientras que en el flujo manual un incidente podía demandar entre 96 y 289 segundos según la complejidad percibida y la carga cognitiva del operador, en el flujo automatizado el rango se comprime a 11--31 segundos, con una varianza casi diez veces menor. Esta compresion de la variabilidad no fue un objetivo explicito del diseñó pero emerge como una consecuencia deseable de la automatización: el sistema no solo es mas rapido en promedio, sino también mas consistente.

La reducción de la intervención humana al 9,5~\%, lejos de ser una limitación, materializa el esquema \emph{human in the loop} buscado: el sistema absorbe la carga rutinaria sin sustituir el juicio humano en los casos que verdaderamente lo demandan. Este equilibrio entre automatización y supervisión preserva la calidad del servicio y libera capacidad operativa significativa. Si se proyecta este porcentaje sobre el volumen trimestral de 3.700 incidentes, la automatización liberaria aproximadamente 180 horas-persona por trimestre, tiempo que el personal podría redirigir a tareas de resolución tecnica o a actividades de mejora continua. La magnitud de esta liberacion ---equivalente a un mes completó de jornada laboral de un operador--- transforma la automatización del registro de un proyecto de optimizacion marginal en una intervención con impacto organizacional medible.

El rigor metodológico adoptado para evitar \emph{data leakage} merece una observación aparte. A diferencia de trabajos que particionan el corpus en conjuntos de entrenamiento y prueba después de realizar exploraciones iniciales ---práctica que puede introducir fuga de información---, el presente estudio aislo completamente la construcción de los componentes críticos del clasificador del corpus de validación de 200 casos. El filtro deterministico se construyó mediante análisis manual de incidentes no incluidos en el corpus, y el prompt fue ajustado en entornos de preproduccion sin exposición a los 200 casos de validación. Solo después de completar el desarrolló se evaluó el desempeño sobre la totalidad del corpus, garantizando que las metricas reportadas reflejan el desempeño genuino del sistema sobre datos no vistos. Esta separación estricta entre datos de construcción y datos de validación constituye una garantía metodológica que no todos los trabajos del área documentan con el mismo nivel de explicitación.

\subsection{Comparacion con la literatura}

Los valores de exactitud y F1 macro obtenidos resultan consistentes con los reportados en la literatura comparable. \textcite{Paramesh2019_it_service_desk} reportan exactitudes cercanas al 87~\% utilizando SVM sobre representaciones TF-IDF, mientras que \textcite{Revina2020_bert_tickets} alcanzan el 91~\% con arquitecturas basadas en BERT. Evaluaciones recientes de LLM en escenarios \emph{few-shot} para clasificación de tickets reportan valores entre el 88~\% y el 94~\% de exactitud dependiendo del idioma y del dominio.

Los resultados del presente trabajo ---92~\% de exactitud global y F1 macro de 0,919--- se ubican dentro del rango superior reportado para idiomas relativamente bien representados en los corpus de entrenamiento de los modelos contemporaneos. Dos comparaciones resultan particularmente informativas. Frente al enfoque SVM+TF-IDF de \textcite{Paramesh2019_it_service_desk} (87~\%), el clasificador hibrido propuesto ofrece una ganancia de 5 puntos porcentuales en exactitud, ganancia que resulta operativamente relevante si se considera que cada error de clasificación implica entre 10 y 15 minutos adicionales de reasignacion. Frente al enfoque BERT de \textcite{Revina2020_bert_tickets} (91~\%), la diferencia es de solo 1 punto porcentual, pero el esquema propuesto no requiere infraestructura de GPU para inferencia local ni un corpus de entrenamiento etiquetado de gran escala, lo que lo hace mas accesible para organizaciones medianas.

La novedad del aporte radica en la validación específica sobre espanol rioplatense aplicado a un dominio de mesa de ayuda de organización mediana, segmento escasamente cubierto por la literatura previa. Diversos trabajos han señalado la menor cobertura del espanol en los corpus de entrenamiento de los LLM contemporaneos, y los resultados del presente trabajo (F1 macro = 0,919) sugieren que, al menos para el subdominio de tickets de soporte, el desempeño en espanol rioplatense es competitivo con el reportado para idiomas de mayor cobertura como el inglés o el alemán.

\subsection{Limitaciones del estudio}

El estudio presenta cinco limitaciones que deben considerarse al interpretar los hallazgos y al planificar replicaciones futuras.

Primero, el tamaño de la muestra (\(n = 200\)) es suficiente para una validación piloto pero insuficiente para sustentar generalizaciones de carácter poblacional sobre el universo de mesas de ayuda en organizaciones medianas. Si bien el intervalo de confianza del 95~\% reportado para la exactitud [87,2~\%, 95,2~\%] no solapa la zona de rechazo, muestras mayores permitirian intervalos mas estrechos y mayor potencia para detectar diferencias finas entre clases.

Segundo, la circunscripcion del corpus a una única organización del sector servicios de la provincia de Mendoza restringe la validez externa de los hallazgos. La distribucion de tipos de incidentes, el lexico organizacional y los patrones de redacción pueden variar significativamente entre organizaciones, sectores y regiones.

Tercero, la dependencia operativa de un servicio externo de inferencia condiciona el rendimiento del sistema. Un eventual deterioro del servicio de Gemini~2.5 Flash, un cambio significativo en las condiciones comerciales del proveedor, o una modificación en el comportamiento del modelo debido a actualizaciones afectaria directamente la operación. Esta dependencia es intrínseca al enfoque basado en LLM externo y debe ser gestionada como un riesgo operativo.

Cuarto, los resultados reflejan el comportamiento del modelo en su versión vigente al momento del estudio (junio de 2026). Las actualizaciones futuras del proveedor pueden modificar tanto el desempeño bruto como la estructura del prompt requerido para alcanzar resultados equivalentes, lo que aconseja la implementación de un régimen de monitoreo continuo de las metricas de clasificación.

Quinto, la comparación se realizó exclusivamente contra el flujo manual (proceso operativo), no contra clasificadores alternativos como TF-IDF+SVM o regresion logistica sobre el mismo corpus. Una evaluación preliminar de un clasificador TF-IDF+SVM sobre un subconjunto de evaluación arrojo una exactitud del 78~\% y un F1 macro de 0,764, lo que confirma que el esquema hibrido propuesto aporta una ganancia sustantiva sobre un \emph{baseline} simple. Investigaciones futuras deberían incluir una comparación sistematica contra un abanico mas amplio de clasificadores sobre el corpus completó.

\subsection{Implicancias practicas}

Desde el punto de vista práctico, los resultados sugieren que organizaciones medianas con volumenes de incidentes comparables al observado en la organización analizada pueden obtener beneficios sustanciales mediante la implementación de soluciones hibridas similares a la propuesta. La reducción de la intervención humana en la etapa de registro libera al personal técnico para concentrarse en la resolución efectiva de los casos, etapa en la cual el juicio experto continua siendo insustituible.

La trazabilidad provista por el registro detallado de cada decision del clasificador habilita dos procesos de mejora continua. Por un lado, permite identificar patrones de error sistematicos y ajustar el filtro deterministico o el prompt del LLM en consecuencia. Por otro lado, la acumulación sostenida de decisiones validadas por los sectores responsables constituye un activo informacional que, una vez alcanzado un volumen adecuado (del orden de 5.000 casos), permitiria entrenar modelos de clasificación específicos del dominio organizacional, con eventuales ventajas en costo, latencia y soberanía sobre los datos.

En el plano económico, el costo operativo del sistema se compone de dos rubros principales: el costo de infraestructura autoalojada (servidor con Docker, estimado en 50--100 USD/mes según el proveedor de nube local) y el costo de inferencia del LLM (aproximadamente 0,02 USD por cada invocacion a Gemini~2.5 Flash en la banda de uso correspondiente). Para un volumen de 3.700 incidentes trimestrales, con un 38~\% requiriendo inferencia externa, el costo trimestral de API se estima en el orden de 28 USD, un valor marginal frente a las 180 horas-persona liberadas.

Un aspecto que excede el alcance técnico de la solución desarrollada, pero que condiciona su efectividad en condiciones reales de operación, es la dimensión organizacional de la adopción. La experiencia recogida durante el trabajo de campo en la mesa de ayuda de referencia indica que los empleados con mayor antiguedad tienden a presentar resistencia frente a cambios de procedimiento, incluso cuando estos cambios ofrecen beneficios demostrables. Esta resistencia no responde a limitaciones tecnicas ni a carencias de la herramienta propuesta, sino a la familiaridad con los procedimientos existentes y a la percepcion ---fundada o no--- de que un nuevo sistema puede resultar complejo o invasivo. En consecuencia, la puesta en produccion de una solución como la aquí presentada se beneficia de una estrategia de despliegue gradual: comenzar con un grupo piloto reducido, extender progresivamente el alcance a sectores adicionales y acompañar cada etapa con instancias de capacitación breves y contextualizadas al dominio de cada área. Un enfoque incremental reduce la friccion inicial, permite ajustar la configuración del sistema a las particularidades de cada sector y facilita la construcción de confianza en la herramienta por parte de los usuarios finales. Esta construcción de confianza constituye una condicion necesaria para que los beneficios técnicos cuantificados en el presente trabajo ---reducción del 89~\% en el tiempo de registro y exactitud del 92~\% en la clasificación--- se traduzcan en mejoras organizacionales efectivas y sostenibles en el tiempo.

\subsection{Reflexiones sobre la generalizacion}

La generalización de los resultados a otras organizaciones requiere atender al menos tres condiciones. En primer lugar, la disponibilidad de un corpus etiquetado de tamaño y calidad similares al utilizado en este trabajo, lo cual implica una inversion inicial en doble etiquetado y validación del acuerdo inter-evaluador. En segundo lugar, la viabilidad tecnica de mantener desplegada localmente la infraestructura de orquestación y persistencia, condicion que descarta a organizaciones sin capacidades operativas minimas para la administración de contenedores Docker. En tercer lugar, la existencia de un marco normativo compatible con el envio de fragmentos descriptivos de incidentes a un servicio externo de inferencia, marco que en el caso argentino se encuentra contemplado por la Ley~25.326 \parencite{Ley25326_2000} bajo determinadas condiciones de minimizacion y pseudonimizacion descritas en el Capitulo~11.

\newpage
